\documentclass[10pt,a4paper]{amsart}
\usepackage[margin=19mm]{geometry}
\usepackage{amsmath,amssymb,mathtools}
\usepackage[scr=rsfs,cal=euler]{mathalfa}
\usepackage{xfrac}
\usepackage[unicode,hidelinks,bookmarks=false]{hyperref}
\newcommand{\ie}{i.e.\ }
\newcommand{\cf}{cf.\ }
\newcommand{\resp}{respectively\ }
\newcommand{\Subsection}{\subsection}
\providecommand{\nobreakdash}{\nobreak\hbox{-}}
\setlength{\emergencystretch}{2em}
\multlinegap0pt
\usepackage{amssymb}
\usepackage{enumerate}
\usepackage[shortlabels]{enumitem}
\usepackage{dsfont}
\usepackage{algorithm}\usepackage{algpseudocode}
\usepackage{mathtools}
\usepackage[normalem]{ulem}
\newtheorem{assumption}{Assumption}
\newtheorem{proposition}{Proposition}
\newcommand{\Dexp}{p}
\newcommand{\F}{\mathcal{F}}
\newcommand{\Imax}{\bar{\imath}}%I_{\max}
\newcommand{\aexp}{\Dexp/(\Dexp-1)}
\newcommand{\f}{\tilde{f}}
\title{Extrapolation-based Direct Search for Nonsmooth Stochastic Zeroth-Order Optimization}
\author{A.~Palmieri \and F.~Rinaldi \and S.~Shashaani}
\date{Source: arXiv:2607.29408v1 (CC BY 4.0)}
\begin{document}
\maketitle

\noindent\emph{Source and licence.} Extrapolation-based Direct Search for Nonsmooth Stochastic Zeroth-Order Optimization. Version arXiv:2607.29408v1 (CC BY 4.0), distributed by arXiv under the Creative Commons Attribution 4.0 International licence, http://creativecommons.org/licenses/by/4.0/, as recorded in the arXiv licence metadata for this exact version and shown on the abstract page https://arxiv.org/abs/2607.29408v1. Full licence notice: \texttt{licenses/arxiv-2607.29408-CC-BY-4.0.txt}.\par\smallskip

\noindent\textit{Editorial scope.} Counted unit: the article's proposition prop:rosenthal-ass33 (source0.tex lines 5297--5422). The article's complete original proof of it is delivered with it (source0.tex lines 5423--5650, one \texttt{proof} environment of 228 lines). No mathematics is written, repaired or re-derived: the statement and the proof are the article's own lines, in the article's own notation, and no retained line is substituted or re-flowed.  Retained uncounted as the source-local context the proof needs: lines 783--812.  Nothing was omitted from any retained span, so the package carries no omission record.  No retained line contains a citation, so the wrapper carries no bibliography and no bibliography entry.  No retained line contains a figure environment, an image-inclusion command or drawing code, so no image is delivered and the package declares no asset dependency.  Editorial formatting only, in addition: an AMS \texttt{amsart} preamble that restates the article's own declarations and macros used by the retained lines, namely the environments \texttt{assumption}, \texttt{proposition} on separate counters, together with the standard packages those lines need.  The retained environments print in this document as Assumption 1, Proposition 1 in order of appearance, whereas the article prints the same items with the numbering of its own full document.  The complete native source of the article as distributed in the arXiv e-print for this exact version is bundled unchanged as \texttt{sources/206481/source0.tex}.
\begin{assumption}[$\Dexp$-tail for single-point errors]\label{ass:quadratic tail}
There exists $\varepsilon_h>0$ such that, for every $\alpha\ge \varepsilon_h$:
\begin{equation}\label{eq:A4}
\begin{aligned}
&\mathbb P\!\left(
|\tilde f_{k}^{(-1)}-f_{k}^{(-1)}|
\ge
\alpha\,\Delta_k^{\Dexp}
\ \middle|\ 
\F_{k-1}
\right)
\le
\frac{\varepsilon_h}{\alpha^{\aexp}},
\\[1mm]
&\mathbb P\!\left(
|\tilde f_{k,m}^{(i)}-f_{k,m}^{(i)}|
\ge
\alpha\,\Delta_k^{\Dexp}
\ \middle|\ 
\F_{k-1}
\right)
\le
\frac{\varepsilon_h}{\alpha^{\aexp}},
\quad
i\in[0:\Imax],\ m\in[1:\bar{m}].
\end{aligned}
\end{equation}


\end{assumption}

\begin{proposition}[Rosenthal batch size for Assumption~\ref{ass:quadratic tail}]
\label{prop:rosenthal-ass33}
Let \(p\in(1,2]\) and set
\[
    r:=\frac{p}{p-1}\ge2.
\]
Let \(\mathcal H\) be a sigma-algebra, let \(X\) be an
\(\mathcal H\)-measurable random point, and let \(\Delta>0\) be an
\(\mathcal H\)-measurable random variable.
Suppose that, conditionally on \(\mathcal H\), the oracle samples
\[
    F(X,\zeta_1),\ldots,F(X,\zeta_W)
\]
are independent and satisfy, for some \(\sigma>0\),
\[
    \mathbb E\!\left[
        F(X,\zeta_j)-f(X)
        \,\middle|\,
        \mathcal H
    \right]
    =0, 
    \qquad 
    \mathbb E\!\left[
        \left|F(X,\zeta_j)-f(X)\right|^r
        \,\middle|\,
        \mathcal H
    \right]
    \le \sigma^r,
\]
for every \(j=1,\ldots,W\), with \(W\) a positive, integer-valued,
\(\mathcal H\)-measurable random variable.\\
Define the sample-average estimator
\[
    \tilde f(X)
    :=
    \frac{1}{W}
    \sum_{j=1}^{W}F(X,\zeta_j).
\]
Then there exists a constant \(C_r>0\), depending only on \(r\), such that,
if
\[
    W
    \ge
    \left(
        \frac{C_r\sigma^r}{\varepsilon_h}
    \right)^{2/r}
    \Delta^{-2p},
\]
then
\[
    \mathbb P\!\left(
        |\tilde f(X)-f(X)|
        \ge
        \alpha\Delta^p
        \,\middle|\,
        \mathcal H
    \right)
    \le
    \frac{\varepsilon_h}{\alpha^{p/(p-1)}}
    \qquad
    \text{for every }\alpha>0.
\]
Consequently, a batch size
 $W=\mathcal O(\Delta^{-2p})$
is sufficient to satisfy Assumption~\ref{ass:quadratic tail}.
% Let $p\in(1,2]$ and set
% \[
%     r:=\frac{p}{p-1}\ge 2.
% \]
% Assume that the stochastic oracle is conditionally unbiased and has uniformly bounded
% conditional $r$-th moment on the set of queried points; namely, for some
% $\sigma_r>0$,
% \[
%     \mathbb E\!\left[
%         F(x,\zeta)-f(x)
%         \,\middle|\,
%         \mathcal F_{k-1}
%     \right]=0,
%     \qquad
%     \mathbb E\!\left[
%         \left|F(x,\zeta)-f(x)\right|^r
%         \,\middle|\,
%         \mathcal F_{k-1}
%     \right]
%     \le
%     \sigma_r^r
% \]
% for every point $x$ queried at iteration $k$.\\
% For each queried point, define the sample-average estimator
% \[
%     \tilde f_{k}^{}(x)
%     :=
%     \frac{1}{w_k^{}}
%     \sum_{j=1}^{w_k^{}}
%     F(x,\zeta_j),
% \]
% where the samples are conditionally independent given $\mathcal F_{k-1}$.\\
% Then there exists a constant $C_r>0$, depending only on $r$, such that if
% \[
%     w_k^{}
%     \ge
%     \left(
%         \frac{C_r\sigma_r^r}{\varepsilon_h}
%     \right)^{2/r}
%     \Delta_{k}^{-2p},
% \]
% then
% \[
% \mathbb P\!\left(
%     |\f_{k}^{}(x)-f(x)|
%     \ge
%     \alpha \Delta_k^p
%     \,\middle|\,
%     \mathcal F_{k-1}
% \right)
% \le
% \frac{\varepsilon_h}{\alpha^{p/(p-1)}}
% \qquad
% \forall \alpha>0.
% \]
% In particular, Assumption~\ref{ass:quadratic tail} is satisfied by considering
% \[
%     w_k^{}=\mathcal O(\Delta_{k}^{-2p}).
% \]
% samples.
\end{proposition}

\begin{proof}
Define
\[
    \xi_j:=F(X,\zeta_j)-f(X),
    \qquad j=1,\ldots,W.
\]
Then
\[
    \tilde f(X)-f(X)
    =
    \frac{1}{W}\sum_{j=1}^{W}\xi_j.
\]
Conditionally on \(\mathcal H\), the random variables
\(\xi_1,\ldots,\xi_W\) are independent and centered, and satisfy
\[
    \mathbb E\!\left[
        |\xi_j|^r
        \,\middle|\,
        \mathcal H
    \right]
    \le\sigma^r.
\]
Since \(r\ge2\), the conditional version of Rosenthal's inequality gives
\[
\begin{aligned}
    \mathbb E\!\left[
        \left|
        \sum_{j=1}^{W}\xi_j
        \right|^r
        \,\middle|\,
        \mathcal H
    \right]
    % &\le
    % C_r\left[
    %     \sum_{j=1}^{w}
    %     \mathbb E\!\left[
    %         |\xi_j|^r
    %         \,\middle|\,
    %         \mathcal H
    %     \right]
    %     +
    %     \left(
    %         \sum_{j=1}^{w}
    %         \mathbb E\!\left[
    %             |\xi_j|^2
    %             \,\middle|\,
    %             \mathcal H
    %         \right]
    %     \right)^{r/2}
    % \right]
    % \\
    % &\le
    \le C_r\sigma^r W^{r/2},
\end{aligned}
\]
where \(C_r>0\) depends only on \(r\).
Therefore,
\[
    \mathbb E\!\left[
        |\tilde f(X)-f(X)|^r
        \,\middle|\,
        \mathcal H
    \right]
    \le
    C_r\sigma^r W^{-r/2}.
\]
By conditional Markov's inequality,
\[
\begin{aligned}
    &\mathbb P\!\left(
        |\tilde f(X)-f(X)|
        \ge
        \alpha\Delta^p
        \,\middle|\,
        \mathcal H
    \right)
    \le
    \frac{
        C_r\sigma^r W^{-r/2}
    }{
        \alpha^r\Delta^{pr}
    }.
\end{aligned}
\]
If
\[
    W
    \ge
    \left(
        \frac{C_r\sigma^r}{\varepsilon_h}
    \right)^{2/r}
    \Delta^{-2p},
\]
then
\[
    C_r\sigma^r W^{-r/2}\Delta^{-pr}
    \le\varepsilon_h.
\]
Hence
\[
    \mathbb P\!\left(
        |\tilde f(X)-f(X)|
        \ge
        \alpha\Delta^p
        \,\middle|\,
        \mathcal H
    \right)
    \le
    \frac{\varepsilon_h}{\alpha^r}.
\]
Since \(r=p/(p-1)\), the claim follows.
% Define
% \[
%     \xi_j(x):=F(x,\zeta_j)-f(x),
%     \qquad j=1,\ldots,w_k.
% \]
% Then
% \[
%     \tilde f_k(x)-f(x)
%     =
%     \frac{1}{w_k}
%     \sum_{j=1}^{w_k}\xi_j(x).
% \]
% Conditionally on $\mathcal F_{k-1}$, the random variables
% $\xi_1(x),\ldots,\xi_{w_k}(x)$ are independent, centered, and satisfy
% \[
%     \mathbb E\!\left[
%         |\xi_j(x)|^r
%         \,\middle|\,
%         \mathcal F_{k-1}
%     \right]
%     \le
%     \sigma_r^r.
% \]
% Since $r\ge2$, the conditional version of Rosenthal's inequality gives
% \[
%     \mathbb E\!\left[
%         \left|
%         \frac{1}{w_k}
%         \sum_{j=1}^{w_k}\xi_j(x)
%         \right|^r
%         \,\middle|\,
%         \mathcal F_{k-1}
%     \right]
%     \le
%     C_r\sigma_r^r w_k^{-r/2},
% \]
% where $C_r>0$ depends only on $r$.\\
% Now apply Markov's inequality to the nonnegative random variable
% $|\tilde f_k(x)-f(x)|^r$. For every $\alpha>0$,
% \[
% \begin{aligned}
% \mathbb P\!\left(
%     |\tilde f_k(x)-f(x)|
%     \ge
%     \alpha\Delta_k^p
%     \,\middle|\,
%     \mathcal F_{k-1}
% \right)
% &=
% \mathbb P\!\left(
%     |\tilde f_k(x)-f(x)|^r
%     \ge
%     \alpha^r\Delta_k^{pr}
%     \,\middle|\,
%     \mathcal F_{k-1}
% \right)
% \\
% &\le
% \frac{
%     \mathbb E\!\left[
%         |\tilde f_k(x)-f(x)|^r
%         \,\middle|\,
%         \mathcal F_{k-1}
%     \right]
% }{
%     \alpha^r\Delta_k^{pr}
% }
% \\
% &\le
% \frac{
%     C_r\sigma_r^r w_k^{-r/2}
% }{
%     \alpha^r\Delta_k^{pr}
% }.
% \end{aligned}
% \]
% If
% \[
%     w_k
%     \ge
%     \left(
%         \frac{C_r\sigma_r^r}{\varepsilon_h}
%     \right)^{2/r}
%     \Delta_k^{-2p},
% \]
% then
% \[
%     C_r\sigma_r^r w_k^{-r/2}\Delta_k^{-pr}
%     \le
%     \varepsilon_h.
% \]
% Therefore,
% \[
% \mathbb P\!\left(
%     |\tilde f_k(x)-f(x)|
%     \ge
%     \alpha\Delta_k^p
%     \,\middle|\,
%     \mathcal F_{k-1}
% \right)
% \le
% \frac{\varepsilon_h}{\alpha^r}.
% \]
% Finally, since $r=p/(p-1)$, this becomes
% \[
% \mathbb P\!\left(
%     |\tilde f_k(x)-f(x)|
%     \ge
%     \alpha\Delta_k^p
%     \,\middle|\,
%     \mathcal F_{k-1}
% \right)
% \le
% \frac{\varepsilon_h}{\alpha^{p/(p-1)}}.
% \]
% Thus Assumption~\ref{ass:quadratic tail} is satisfied.
\end{proof}

\end{document}
